\section{What the quasi-Riemann hypothesis can contribute}
\label{sec:qrh-contribution}

The two supplied quasi-Riemann-hypothesis papers have different endpoints.\cite{qrh-september,qrh-october}
The paper of 30 September 2026 states zero exclusion in
$\operatorname{Re}s>7/8$ for every finite-order Hecke character of
$\mathbb Q(\sqrt{-3})$, and hence for every Dirichlet character.
The shorter paper of 5 October states the corresponding $11/12$ result.
The latter reduces its conclusion to a particularly concrete mean-square
estimate (its Proposition~3.1). The reductions below can be checked
independently of the substantial cubic-theta calculation that supplies
that estimate. We do not treat that calculation as an elementary lemma
or silently include it among the hypotheses already proved here.

For use below, write $\mathcal Z(\vartheta)$ for the following precise
assertion, where $1/2\leq\vartheta<1$:
\[
 L(s,\chi)\ne0\quad(\operatorname{Re}s>\vartheta)
\]
for every Dirichlet character $\chi$, with the pole of a principal
character at $1$ allowed. A statement about each character with its
own unspecified boundary is weaker than $\mathcal Z(\vartheta)$.

\subsection{A direct Mellin-transform criterion}

For a number field $K$, a finite-order Hecke character may be
viewed as a character of a ray class group: it is a homomorphism from
the group of fractional ideals prime to a fixed integral modulus
$\mathfrak f$ to the roots of unity, trivial on the principal ideals
$(\alpha)$ with $\alpha\equiv1\pmod{\mathfrak f}$ and with
$\alpha$ positive at the real places prescribed by the modulus.
Let $\nu$ be such a character, and let $S$ be a finite set of prime
ideals containing the primes dividing its conductor. The notation
$(\mathfrak a,S)=1$ means that no member of $S$ divides
$\mathfrak a$. Let $\mu_K$ be the ideal M\"obius function: it is zero on ideals
divisible by a prime square and equals $(-1)^r$ on a product of $r$
distinct primes. The norm of an integral ideal $\mathfrak a$ is
$N\mathfrak a=|\mathcal O_K/\mathfrak a|$. Put
\[
 L^S(s,\nu)=\prod_{\mathfrak p\notin S}
       (1-\nu(\mathfrak p)N\mathfrak p^{-s})^{-1}
       \qquad(\operatorname{Re}s>1).
\]
Assume that the number of ideals of norm at most $t$ is $O_K(t)$,
and that $L^S$ has its usual meromorphic continuation to
$\operatorname{Re}s>0$. These standard analytic inputs will be kept
explicit in this subsection.

\begin{proposition}[Compactly supported smoothing]
\label{prop:qrh-mellin}
Fix $0<\alpha<1$. Suppose that, for every
$W\in C_c^\infty(0,\infty)$ and every $\epsilon>0$,
\[
 A_W(X):=\sum_{(\mathfrak a,S)=1}
     \mu_K(\mathfrak a)\nu(\mathfrak a)
     W(N\mathfrak a/X)
       =O_{K,\nu,S,W,\epsilon}(X^{\alpha+\epsilon}).
\]
Then $L(s,\nu)$ has no zero in $\operatorname{Re}s>\alpha$.
No uniformity in $\nu$ of the implied constant is required.
\end{proposition}
\begin{proof}
If $\rho$ were such a zero, choose a nonzero nonnegative function
$\phi\in C_c^\infty(1,2)$ and set
$W(y)=y^{-\rho}\phi(y)$, where $y^{-\rho}=\exp(-\rho\log y)$.
The Mellin transform
\[
 \widehat W(s)=\int_0^\infty W(y)y^s\,\frac{dy}{y}
\]
is entire, and $\widehat W(\rho)=\int\phi(y)\,dy/y>0$.
On each compact subset of $\operatorname{Re}s>\alpha$, choose
$\epsilon$ smaller than the distance of that compact set from the
line $\operatorname{Re}s=\alpha$. The assumed bound, including the
same bound after multiplying the integrand by any fixed power of
$\log X$, proves local uniform convergence and holomorphy of
\[
 M_W(s)=\int_0^\infty A_W(X)X^{-s}\,\frac{dX}{X}.
\]
There is no issue at zero: compact support makes $A_W(X)=0$ for all
sufficiently small $X$. For $\operatorname{Re}s>1$, absolute
convergence permits termwise integration and the substitution
$y=N\mathfrak a/X$ gives
\[
 M_W(s)=\widehat W(s)
       \sum_{(\mathfrak a,S)=1}
        \frac{\mu_K(\mathfrak a)\nu(\mathfrak a)}{N\mathfrak a^s}
       =\frac{\widehat W(s)}{L^S(s,\nu)}.
\]
Consequently the meromorphic functions satisfy
$L^S(s,\nu)M_W(s)=\widehat W(s)$ throughout the connected
half-plane $\operatorname{Re}s>\alpha$, by the identity theorem.
At a zero $\rho$, both $L^S$ and $M_W$ are holomorphic, so evaluation
gives $0=\widehat W(\rho)>0$. Finally, the finitely many factors
$1-\nu(\mathfrak p)N\mathfrak p^{-s}$ have no zeros in
$\operatorname{Re}s>0$, because their second term has absolute
value strictly less than one. Thus deleting them does not change
the zeros under consideration.
\end{proof}

There is an even simpler criterion when an exponentially smoothed
sum can be estimated directly. Define
\[
 F_\chi(X)=\sum_{n\ge1}\mu(n)\chi(n)e^{-n/X}.
\]
If $F_\chi(X)=O_{\chi,\epsilon}(X^{\alpha+\epsilon})$ for every
$\epsilon>0$, then
\[
 \int_0^\infty F_\chi(X)X^{-s}\,\frac{dX}{X}
       =\frac{\Gamma(s)}{L(s,\chi)}
       \qquad(\operatorname{Re}s>1)
\]
extends holomorphically to $\operatorname{Re}s>\alpha$.
At zero, the absolute sum is bounded by
$e^{-1/X}/(1-e^{-1/X})$; at infinity use the assumed estimate.
The preceding proof therefore applies with the single factor
$\Gamma(s)$, which has no zeros. Thus it suffices to estimate one
fixed smooth function; a collection of auxiliary functions is not
inherent in the analytic implication. This is a simplification of
the criterion, not a proof of the required cancellation estimate.

\subsection{The exact cost of amplification by powers of primes}

The following lemma isolates the numerical gain and the error that
must be paid. It applies to ideals as well as to positive integers.
We state it for $K=\mathbb Q(\sqrt{-3})$, where every ideal prime
to $3$ has a unique generator congruent to $1$ modulo $3$.
Write $\mathcal O=\mathbb Z[\omega]$, with
$\omega=(-1+\sqrt{-3})/2$. The six units give the six invertible
classes modulo $3\mathcal O$, so the asserted uniqueness follows
from the fact that $\mathcal O$ is a Euclidean domain.

\begin{lemma}[Power amplification]
\label{lem:qrh-amplification}
Let $k\ge2$, and let $\chi_{\mathfrak a}(u)$ be multiplicative
characters in $u$, of order dividing $k$ away from
$\mathfrak a$, extended by zero when $(u,\mathfrak a)\ne1$.
For a fixed bounded smooth function $W$ supported in a compact
subinterval of $(0,\infty)$, put
\[
 A_u(D)=\sum_{(\mathfrak a,S)=1}
 \mu_K(\mathfrak a)\nu(\mathfrak a)
 \chi_{\mathfrak a}(u)W(N\mathfrak a/D).
\]
Suppose $2\le Y=H^{1/k}\le D$, and that
\[
 \sum_{0<Nu\le H}|A_u(D)|^2\le B(D,H).
\]
If there are $\gg Y/\log Y$ prime ideals outside $S$ with norms
in $(Y/2,Y]$, then
\[
 |A_1(D)|^2\ll
       B(D,H)H^{-1/k}\log(2H)+D^2H^{-2/k}.
\]
The constants can depend on $K,S,W$ and on the constant in the
prime count, but not on $D$ or $H$.
\end{lemma}
\begin{proof}
Let $p$ be the primary generator of one of the indicated prime
ideals. The convention at nonunits gives exactly
\[
 \chi_{\mathfrak a}(p^k)=1_{p\nmid\mathfrak a}.
\]
Thus $A_1(D)-A_{p^k}(D)$ is a sum over ideals divisible by $p$.
There are $O_W(D/Np)$ such ideals, because the remaining ideal
has norm at most a fixed multiple of $D/Np$; the condition
$Y\le D$ absorbs the constant in ideal counting. Hence
\[
 |A_1(D)-A_{p^k}(D)|\ll_W D/Y.
\]
Distinct primary primes have distinct $k$th powers. If there
are $J\gg Y/\log Y$ of them, the inequality
$|z+w|^2\le2|z|^2+2|w|^2$ gives
\[
 J|A_1(D)|^2\le
 2\sum_p|A_{p^k}(D)|^2+O_W(JD^2/Y^2).
\]
Their norms are at most $H$, and the first sum is bounded by
$B(D,H)$. Divide by $J$.
\end{proof}

Suppose more generally that
$B(D,H)\ll_\epsilon D^{a+\epsilon}H^b$ and $H=D^h$.
The exponent furnished by this argument is
\begin{equation}\label{eq:qrh-general-exponent}
 \max\left\{\frac a2+\frac h2\left(b-\frac1k\right),
                    1-\frac hk\right\},
\end{equation}
up to an arbitrary positive loss. In the October paper,
$k=6$, $a=b=1$, and the proved range used for the mean square
is $H=D^{1+\eta}$, $0<\eta\le1/10$. Therefore the first
entry is $11/12+5\eta/12$, the second is
$5/6-\eta/6$, and letting $\eta$ tend to zero gives $11/12$.
This verifies the final deduction from Proposition~3.1 without
any claim about its difficult proof.

Even if the bound $B(D,H)\ll_\epsilon D^{1+\epsilon}H$
were known for every $H=D^h$, the best exponent obtainable
from \emph{this estimate together with this absolute bound for
prime deletion} would be $k/(k+1)$. Indeed the two entries in
\eqref{eq:qrh-general-exponent} are equal at
$h=k/(k+1)$, and one increases while the other decreases.
For $k=6$ this gives $6/7$, still greater than $1/2$.
This is a limitation of these two inequalities, not an
impossibility theorem for character amplification.

One can retain more arithmetic than the absolute deletion
bound. For $p\notin S$, multiplicativity of the M\"obius function
gives the exact identity
\begin{equation}\label{eq:qrh-prime-removal}
 A_1(D)=A_{p^k}(D)-\nu(p)A_{p^k}(D/Np).
\end{equation}
To prove it, split the defining sum according to $p\mid\mathfrak a$;
write $\mathfrak a=p\mathfrak b$ in the latter part, use
$\mu_K(p\mathfrak b)=-\mu_K(\mathfrak b)$ when $p\nmid\mathfrak b$,
and use zero otherwise. An a priori bound for $A_1$ can be used
in this identity to control prime deletion uniformly. Section~\ref{sec:qrh-bootstrap}
proves that improvement and states the additional mean-square estimate
it would require. A bound whose constant is allowed to depend on the
varying prime cannot be inserted into the amplification without a
separate uniformity argument.

For completeness, the positivity step alone cannot force
an exponent below $11/12$ in the range $H=D$. Consider abstract
complex numbers with $J$ distinguished entries, all equal to
$D/\sqrt J$, and all other entries zero. Their sum of squares
is exactly $D^2$, while the common distinguished value has
size $D^{11/12}$ when $J=D^{1/6}$.
This is an exact example for the numerical inequalities, not
an example consisting of M\"obius sums and not a counterexample
to either submitted theorem.

\subsection{Passing from the Eisenstein field to Dirichlet characters}

\begin{proposition}[Quadratic base change at the analytic endpoint]
\label{prop:qrh-basechange}
Let $0<\alpha<1$. If every finite-order Hecke $L$-function of
$K=\mathbb Q(\sqrt{-3})$ has no zero in
$\operatorname{Re}s>\alpha$, with principal poles allowed,
then the same conclusion holds for every Dirichlet $L$-function.
\end{proposition}
\begin{proof}
Let $\epsilon$ be the nontrivial character modulo $3$.
For a Dirichlet character $\chi$, pull back its ideal character
along the norm from $K$ to $\mathbb Q$. At primes where all
characters are unramified, the Euler factors give
\[
 L_K(s,\chi\circ N)=L(s,\chi)L(s,\chi\epsilon).
\]
If a rational prime splits, its two prime ideals have norm $p$
and $\epsilon(p)=1$, so both sides have factor
$(1-\chi(p)p^{-s})^{-2}$. If it is inert, its prime ideal
has norm $p^2$ and $\epsilon(p)=-1$, so both sides have
factor $(1-\chi(p)^2p^{-2s})^{-1}$.
At the finitely many excluded primes insert the usual Euler
factors; they are nonzero for $\operatorname{Re}s>0$.
Equality in the half-plane of absolute convergence extends
meromorphically.

Away from $1$, a zero of either Dirichlet factor would give
a zero of the product, because neither factor has a pole
there. At $1$, a cancellation could occur only when one of
the factors is principal, so the other is induced by
$\epsilon$. But
\[
 L(1,\epsilon)=\sum_{j=0}^\infty
      \left(\frac1{3j+1}-\frac1{3j+2}\right)>0.
\]
The series converges and every displayed summand is positive.
Imprimitive Euler factors at $1$ are also nonzero. Thus a
pole cannot conceal a zero, completing the proof.
\end{proof}

\subsection{A useful uniform consequence: logarithmic derivatives}

We now give a consequence of $\mathcal Z(\vartheta)$ that
really does shorten one of the auxiliary-prime arguments in
the submitted proof. The hypothesis remains explicit: the
following proof is independent of how $\mathcal Z(\vartheta)$
is obtained.

\begin{lemma}[A logarithm on a disk]
\label{lem:qrh-log-disk}
Let $f$ be holomorphic on $|z-z_0|<R$, and suppose
$\operatorname{Re}f\le M$ there. For $0<r<R$,
\[
 \sup_{|z-z_0|\le r}|f(z)-f(z_0)|
 \le \frac{2r}{R-r}\bigl(M-\operatorname{Re}f(z_0)\bigr).
\]
\end{lemma}
\begin{proof}
Put $g=f-f(z_0)$ and $A=M-\operatorname{Re}f(z_0)\ge0$.
If $A=0$, the maximum principle for the real part makes
$g=0$. If $A>0$, the function $g/(2A-g)$ maps the disk
into the closed unit disk, since $\operatorname{Re}g\le A$,
and vanishes at $z_0$. Schwarz's lemma after rescaling gives
$|g/(2A-g)|\le |z-z_0|/R$. Rearrangement gives the stated
bound. Applying the argument first on smaller closed disks
and taking a limit avoids assumptions at the boundary.
\end{proof}

\begin{proposition}[Uniform logarithmic derivative]
\label{prop:qrh-log-derivative}
Assume $\mathcal Z(\vartheta)$. For every fixed
$\sigma>\vartheta$, there is a constant $C_{\vartheta,\sigma}$
such that, for every nonprincipal Dirichlet character modulo
$q$ and all real $t$,
\[
 \left|\frac{L'}{L}(\sigma+it,\chi)\right|
 \le C_{\vartheta,\sigma}\log\bigl(q(|t|+3)\bigr).
\]
The same estimate holds at every point in any fixed closed
vertical strip whose real parts lie in
$(\vartheta,\infty)$. In particular
$|L'(1,\chi)/L(1,\chi)|\ll_\vartheta\log(q+3)$.
For a principal character, subtract its simple-pole term
before making the corresponding assertion.
\end{proposition}
\begin{proof}
It is enough first to consider primitive nonprincipal
characters and $\vartheta<\sigma\le2$. The partial sums
$A(u)=\sum_{n\le u}\chi(n)$ have absolute value at most $q$:
a complete period sums to zero. Partial summation therefore
gives, for $\operatorname{Re}s>0$,
\[
 L(s,\chi)=s\int_1^\infty A(u)u^{-s-1}\,du,
 \qquad |L(s,\chi)|\le\frac{q|s|}{\operatorname{Re}s}.
\]
This also proves holomorphy there by locally uniform
convergence. Choose radii
\[
 2-\sigma<r<R<2-\vartheta
\]
depending only on $\vartheta,\sigma$. On the disk of radius
$R$ centered at $2+it$, the zero-free hypothesis gives a
holomorphic logarithm $f=\log L$. Choose its value at the
center by the absolutely convergent Euler-product logarithm;
that value is bounded absolutely, uniformly in $q,t$.
The preceding bound gives
$\operatorname{Re}f\le\log(q(|t|+3))+C_{\vartheta,\sigma}$
on the larger disk. Lemma~\ref{lem:qrh-log-disk} bounds $f$
on the disk of radius $r$ by the same quantity.
A further fixed smaller circle around $\sigma+it$ lies
inside that disk. Cauchy's derivative formula bounds
$f'(\sigma+it)=L'/L(\sigma+it,\chi)$ as asserted.
The argument is uniform as $\sigma$ varies in a compact
subinterval of $(\vartheta,2]$. For larger real parts the
Euler product gives the bound directly.

If $\chi$ is induced from $\chi^*$, then
\[
 \frac{L'}L(s,\chi)=\frac{L'}L(s,\chi^*)+
 \sum_{\substack{p\mid q\\p\nmid q^*}}
       \frac{\chi^*(p)\log p}{p^s-\chi^*(p)}.
\]
For $\operatorname{Re}s\ge\sigma>0$, the absolute value of
the sum is at most
$(2^\sigma-1)^{-1}\sum_{p\mid q}\log p
\le(2^\sigma-1)^{-1}\log q$.

For the primitive principal character use
$F(s)=(s-1)\zeta(s)$ instead of $L(s,\chi)$. It is
holomorphic and zero-free in the relevant disks, including
at $1$. The elementary continuation formula
\[
 \zeta(s)=\frac{s}{s-1}
          -s\int_1^\infty\{u\}u^{-s-1}\,du
\]
bounds $F$ polynomially in $|t|+3$ on those disks.
At their centers a logarithm of $F$ is
$\log(1+it)+\log\zeta(2+it)$, of size
$O(\log(|t|+3))$. The same argument bounds $F'/F$;
finally $\zeta'/\zeta=F'/F-1/(s-1)$.
\end{proof}

\subsection{Primes in progressions without an exceptional-zero case}

Let $\Lambda(n)$ equal $\log p$ when $n=p^j$ is a prime
power, and zero otherwise. Choose a fixed nonnegative
$\Phi\in C_c^\infty(1,2)$ with $\int\Phi(y)\,dy>0$.

\begin{proposition}[A smooth prime-counting estimate]
\label{prop:qrh-smooth-primes}
Assume $\mathcal Z(\vartheta)$ and fix
$\vartheta<\sigma<1$. Uniformly in $q\ge1$, $(a,q)=1$
and $X\ge2$,
\[
 \sum_{n\equiv a\ (q)}\Lambda(n)\Phi(n/X)
 =\frac{X}{\varphi(q)}\int_0^\infty\Phi(y)\,dy
   +O_{\vartheta,\sigma,\Phi}
      \bigl(X^\sigma\log(q+3)\bigr).
\]
Consequently, for every $B>1/(1-\vartheta)$ there is
$C_B$ such that every reduced class modulo $q$ contains
a prime in $(X,2X)$ whenever $X\ge C_B(q+1)^B$.
\end{proposition}
\begin{proof}
For $c>1$, absolute convergence of the Euler product and
Mellin inversion give
\[
 \sum_n\Lambda(n)\chi(n)\Phi(n/X)
 =\frac1{2\pi i}\int_{(c)}
    -\frac{L'}L(s,\chi)\widehat\Phi(s)X^s\,ds.
\]
Mellin inversion here is ordinary Fourier inversion after
the substitution $y=e^u$. Repeated integration by parts in
$u$ gives, for every integer $A\ge0$,
$|\widehat\Phi(v+it)|\ll_{A,\Phi}(1+|t|)^{-A}$,
uniformly for $v$ in a fixed compact interval.
Proposition~\ref{prop:qrh-log-derivative} therefore permits
moving the line from $c$ to $\sigma$. On the horizontal
segments the integral tends to zero; the vertical integral
is absolutely bounded by
$O(X^\sigma\log(q+3))$, since
$\log(q(|t|+3))\le\log(q+3)+\log(|t|+3)$.
The only pole crossed is at $1$ for a principal character,
with residue $X\widehat\Phi(1)$.
Character orthogonality gives the progression formula.
There are $\varphi(q)$ terms in that average and the factor
$1/\varphi(q)$ cancels their number, so the error does not
acquire an extra factor of $\varphi(q)$.

Terms with $n=p^j$, $j\ge2$, contribute
$O_\Phi(X^{1/2}(\log X)^2)$: for each
$2\le j\le\log(2X)/\log2$ there are at most
$(2X)^{1/j}\le(2X)^{1/2}$ possibilities for $p$, and
$\log p\le\log(2X)$. Since $\sigma>1/2$, this is $O_\sigma(X^\sigma)$,
although it is useful to retain its smaller exponent when
comparing it with the main term.
Choose $\sigma$ with $\vartheta<\sigma<1-1/B$.
For $X\ge C_B(q+1)^B$, the ratios of both errors to
$X/\varphi(q)$ tend to zero uniformly as $C_B$ tends
to infinity: use $\varphi(q)\le q$,
$qX^{\sigma-1}\log(q+3)\ll
 C_B^{\sigma-1}(q+1)^{1-B(1-\sigma)}\log(q+3)$,
and the analogous bound for prime powers.
The weighted sum over actual primes is consequently
positive, proving existence.
\end{proof}

\begin{corollary}[The progression input in the submitted proof]
\label{cor:qrh-ProgAt}
Assume $\mathcal Z(\vartheta)$. There is a constant $C=C_\vartheta\ge1$ such that, for every positive integer $M$,
every integer $a$ coprime to $M$, and every
$X\ge\max(3,M^C)$, there is a prime $p$ satisfying
\[
 X<p\le X^2,\qquad p\equiv a\pmod M,
 \qquad
 \left|\frac{L'(1,\chi_p)}{L(1,\chi_p)}\right|
      \le C\log p,
\]
where $\chi_p(n)=(n/p)$ is the quadratic character modulo
$p$. For the fixed boundary $\vartheta=7/8$, the constant is absolute.
This supplies the ordinary progression input used in the
submitted argument.
\end{corollary}
\begin{proof}
For $M\ge2$, enlarge $C$ so that $M^C$ dominates
$C_B(M+1)^B$ in Proposition~\ref{prop:qrh-smooth-primes}.
The prime then lies in $(X,2X)\subset(X,X^2]$.
Proposition~\ref{prop:qrh-log-derivative} at $s=1$ gives
the logarithmic derivative bound for every such prime,
not merely for a selected subfamily. For $M=1$,
Bertrand's postulate supplies a prime in $(X,2X)$ for
$X\ge3$, and the same bound applies. One may also cover
bounded $X$ by finitely many primes and increase $C$.
\end{proof}

\begin{corollary}[The two auxiliary quadratic fields]
\label{cor:qrh-auxiliary-pair}
Assume $\mathcal Z(\vartheta)$ and fix $\epsilon>0$.
For all sufficiently large $V$ there are distinct primes
$q,d$ such that
\[
 q\equiv3\pmod8,\quad d\equiv7\pmod{16},\quad q\mid d+1,
 \quad p\mid d+1\quad
  \text{for every odd prime }p\le V^{\epsilon/5},
\]
\[
 \tfrac12V^{\epsilon/5}<\log q\le V^{\epsilon/5},\qquad
 \tfrac12V^{\epsilon/5}(\log V)^2<\log d
                         \le V^{\epsilon/5}(\log V)^2,
\]
and
\[
 \left|\frac{L'(1,\chi_d)}{L(1,\chi_d)}\right|
       \le V^\epsilon.
\]
In particular $(-d/p)=1$ for every indicated odd prime $p$.
Since $d\equiv3\pmod4$, quadratic reciprocity identifies
$\chi_d(n)=(n/d)$ with the primitive quadratic character
$\bigl(\frac{-d}{n}\bigr)$ of the field
$\mathbb Q(\sqrt{-d})$. For odd $n$ the reciprocity signs
cancel against $(-1/n)$; the supplementary law at $2$ gives
the same equality for even $n$, and multiplicativity completes
the verification.
\end{corollary}
\begin{proof}
Put $Y=V^{\epsilon/5}$. Apply
Proposition~\ref{prop:qrh-smooth-primes} to the class $3$
modulo $8$ at $X=e^{3Y/4}$. The resulting prime $q$
has the required logarithmic bounds and satisfies $q>Y$
for large $V$.
Let $P=\prod_{3\le p\le Y}p$ and $M=16qP$.
The Chinese remainder theorem gives a reduced class
$a\pmod M$ with $a\equiv7\pmod{16}$ and
$a\equiv-1\pmod{qP}$.
We only need the elementary bound
$\log P\le Y\log Y$: there are at most $Y$ factors,
each at most $Y$. Thus $\log M=O_\epsilon(Y\log V)$.
Take $X=e^{3Y(\log V)^2/4}$. For every fixed $B$,
$X\ge C_B(M+1)^B$ once $V$ is sufficiently large.
A prime $d\in(X,2X)$ in the class $a$ has all the
required size and congruence conditions. It is distinct
from $q$ because $d$ is much larger. Moreover
$-d\equiv1\pmod p$ proves $(-d/p)=1$.
Finally Proposition~\ref{prop:qrh-log-derivative} bounds
the logarithmic derivative by
$O_\vartheta(\log d)=O_\vartheta(Y(\log V)^2)$,
which is less than $Y^5=V^\epsilon$ eventually.
\end{proof}

These corollaries remove an exceptional-zero selection
argument from the ordinary Dirichlet part of the proof.\cite{alpo-prime-supply}
They do not replace the distinct assertion about primes
represented by
\[
 (1+Mu)^2+d(Mv)^2.
\]
That assertion is a prime theorem for principal ray classes
in the varying imaginary quadratic field
$\mathbb Q(\sqrt{-d})$. A Hecke character of this field
need not be a norm pullback of a Dirichlet character.
Neither $\mathcal Z(\vartheta)$ nor the supplied theorem
for the one fixed field $\mathbb Q(\sqrt{-3})$ supplies
that varying-field theorem. Retaining this distinction is
essential for a valid simplification.

\subsection{Why these estimates do not supply the remaining arithmetic argument}

In the final scalar estimate of the submitted argument,\cite{alpo-final-scalar}
$A\le2xb$ is an upper bound for an auxiliary height or
weight parameter. It is not a bound for a M\"obius sum.
After division by the common positive factor, its contribution
includes $A/(mb)$, while another contribution is proportional
to $m$. Thus the elementary balance is $x/m+m$, whose
minimum is $2\sqrt{x}$. A proof of
$\mathcal Z(7/8)$ does not change $A$ into
$O(bx^{7/8})$.

The RH subcase separately requires square-root bounds for
three selected quadratic characters. The full Dirichlet-GRH
construction requires them for a finite family of twists of
fixed characters by the selected quadratic characters, uniformly
on the chosen interval of scales. A uniform ordinary bound
$O(x^{7/8+\epsilon})$ is weaker and cannot be substituted
directly. However, selection by averaging can use the fixed
zero-free half-plane in a different way; the continuation below
gives a conditional proof of that selection, including the fixed
twists. The estimates for the theta lifts, the local divisibility
and operator norms, and the Galois cohomology still remain
separate arithmetic inputs.

\subsection{Two checks on proposed analytic shortcuts}

The catalogue is useful for locating exact reductions and
for recording unsuccessful estimates.\cite{qrh-catalogue} Its statements that
a component is conditional, source-relative, or only checked
as an algebraic inequality have mathematical force: they
cannot be dropped when using the component. The following
obstruction, corresponding to its entry G280, admits a
short complete proof.

\begin{proposition}[Finite Dirichlet polynomials and absolute convergence]
\label{prop:qrh-absolute-obstruction}
Let $\chi$ be a Dirichlet character and let
$B(s)=\sum_{d\le D}b_d d^{-s}$ be a finite Dirichlet
polynomial. Expand
$1-L(s,\chi)B(s)=\sum_{n\ge1}r_n n^{-s}$ for
$\operatorname{Re}s>1$. For real $\sigma\le1$, the
inequality $\sum_n|r_n|n^{-\sigma}<1$ is impossible.
\end{proposition}
\begin{proof}
If $b_1=0$, then $r_1=1$. If $b_1\ne0$, every prime
$p>D$ coprime to the modulus has
$r_p=-b_1\chi(p)$, so $|r_p|=|b_1|$.
The series $\sum_p1/p$ diverges. To recall its elementary
proof, convergence would bound
$\prod_{p\le y}(1-p^{-1})^{-1}$ uniformly in $y$;
expanding that product gives at least
$\sum_{n\le y}1/n$, a contradiction. Removing finitely
many primes does not alter divergence. Since
$p^{-\sigma}\ge p^{-1}$ for $\sigma\le1$, the
absolute residual series is infinite.
\end{proof}

A different technical issue is whether a Mellin multiplier
could vanish at the very zero one wants to exclude.
Catalogue entry G276 suggests the following useful general
fact, which can be proved without asymptotic integration.

\begin{proposition}[A finite family of Mellin transforms]
\label{prop:qrh-mellin-family}
Let $a>0$ and let $P(y)=\sum_{j=0}^d c_jy^j$ be a real
polynomial of degree $d\ge1$. For
$\operatorname{Re}s>0$ put
\[
 F(s)=\int_0^\infty y^{s-1}e^{-ay-iP(y)}\,dy.
\]
The $d$ numbers $F(s),F(s+1),\ldots,F(s+d-1)$ never
vanish simultaneously.
\end{proposition}
\begin{proof}
Exponential decay and local integrability at zero justify
holomorphy and differentiation under the integral on every
compact subset of $\operatorname{Re}s>0$.
Integration of the derivative of
$y^s e^{-ay-iP(y)}$ gives
\[
 sF(s)=aF(s+1)+i\sum_{j=1}^d j c_jF(s+j).
\]
For $d\ge2$, the coefficient $idc_d$ is nonzero.
If the first $d$ consecutive values vanished, this relation,
then its translates by positive integers, would make every
$F(s+n)$ vanish. For $d=1$ the same conclusion uses
$a+ic_1\ne0$.

Now the finite complex measure
$y^{s-1}e^{-ay-iP(y)}\,dy$ has Laplace transform
\[
 H(z)=\int_0^\infty y^{s-1}e^{-ay-iP(y)-zy}\,dy,
       \qquad\operatorname{Re}z>-a.
\]
All its derivatives at zero would vanish. Holomorphy
would give $H\equiv0$. Push the measure forward under
$u=e^{-y}$ to $[0,1]$. Its moments are $H(n)=0$ for
all integers $n\ge0$. Polynomial approximation of
continuous functions on $[0,1]$ and finite total variation
then show that this measure is zero. Its original density
is nonzero everywhere on $(0,\infty)$, a contradiction.
\end{proof}

This last proposition removes a possible cancellation by
a fixed Mellin multiplier. It supplies no bound for the
arithmetic sums multiplying that factor. In particular,
it cannot turn a bound with exponent $7/8$ into one with
exponent $1/2$.

\subsection{An exact cancellation in the double Poisson argument}

The reindexing in Section 7.3 of the October paper contains a useful
finite identity. Its strength comes from performing the cancellation
before replacing any summand by its absolute value. The following
weighted form is slightly more general and has a short proof.

\begin{proposition}[A squarefree factorization identity]
\label{prop:qrh-factorization}
Let $f$ be squarefree in a unique factorization domain, or let
$f$ be a squarefree integral ideal of a number field. Let $k\ne0$
be an element, or a nonzero integral ideal, respectively.
Assign an arbitrary complex number $c(p)$ to each prime dividing
$f$, and extend $c$ multiplicatively to its squarefree divisors.
Then
\[
 \sum_{Cew=f}\mu(e)\mu(w)c(e)
    1_{e\mid k}\,1_{(w,k)=1}
           \sum_{d\mid(C,k)}c(d)
       =1_{f\mid k}.
\]
Factorizations count associates only once in the element formulation.
Since $f$ is squarefree, the three factors $C,e,w$ are automatically
pairwise coprime.
\end{proposition}
\begin{proof}
Both sides factor over the primes dividing $f$. Fix such a prime
$p$, and put $I=1_{p\mid k}$. Placing $p$ in $C$ contributes
$1+c(p)I$, placing it in $e$ contributes $-c(p)I$, and placing
it in $w$ contributes $-(1-I)$. Their sum is $I$. Multiplication
over $p\mid f$ gives the result. The empty product handles $f=1$.
\end{proof}

For $c(p)=1$ this is the cancellation used after the two Poisson
summations: it forces the new auxiliary divisor to divide the new
frequency. It is valid at shared primes, including the terms where
that frequency is divisible by a prime in the original exclusion.
Dropping those terms would change the identity. The weighted form
shows that bounded prime weights can sometimes be retained through
this particular step; a theorem about the complete surrounding
mean square would still have to pay all other terms and the
resulting uniform constants.


